\chapterimage{figs/chapterhead}
\chapter*{Preface}
\addcontentsline{toc}{chapter}{\textcolor{ocre}{Preface}}

This book presents the \textit{GenESyS} simulator as both a tool for practical use and a development platform. Its purpose is twofold: it serves the reader who wants to install the system, understand how it is used, build models, execute them, and interpret the results; it also serves the developer who wants to understand the internal architecture and extend the simulator with new components, data definitions, language functions, tools, and applications.

Unlike a broad text on systems modeling and simulation, this volume is deliberately centered on software. General concepts appear only when they are needed to explain how \textit{GenESyS} works, how it is organized, or how its resources should be used. The emphasis remains on concrete repository-backed material: operational descriptions, documented commands, architectural decisions, and reproducible procedures.

The book is organized to support progressive reading. The early chapters focus on direct use: installation, first execution, model creation, simulation, and result inspection. The middle chapters explain the user-facing language, shell, worker boundary, and standalone tools. The later chapters move into architecture, implementation, testing, packaging, and governance.

The front matter below makes the reading strategy and evidence model explicit. It records the repository snapshot used for this revision, the main paths through the book, the current safety and scientific boundaries, and the conventions used throughout the manual.

\section*{Repository Revision and Evidence}
\addcontentsline{toc}{section}{Repository Revision and Evidence}
\label{sec:front-matter-repository-revision}

This manual revision is anchored to the current repository snapshot used for this update:

\begin{itemize}
\item Repository: \path{rlcancian/Genesys-Simulator}
\item Branch: \path{WorkInProgress}
\item Commit: \path{f3214dfbde92b28a292416f2abbb033af023d25d}
\item Manual revision date: July 23, 2026
\item Build environment: Ubuntu 24.04.4 LTS
\item Compiler: g++ 13.3.0
\item CMake: 3.28.3
\item Qt: 6.4.2
\end{itemize}

All source-level claims in this book should be read against that snapshot unless a later chapter explicitly states a newer validated commit.

\section*{Reading Paths}
\addcontentsline{toc}{section}{Reading Paths}
\label{sec:front-matter-reading-paths}

The manual is written for several entry points:

\begin{itemize}
\item first contact;
\item GUI modeling;
\item terminal execution;
\item component development;
\item architecture;
\item tests;
\item scientific simulation;
\item biological extensions.
\end{itemize}

Figure~\ref{fig:manual-reading-map} shows the recommended progression through the book.

% FIGURE-SPEC-BEGIN
% ID: fig:manual-reading-map
% Source asset: figs/generated/manual-reading-map
% Caption: Recommended progression through the manual from first contact to advanced topics.
% Accessibility: A structured map that highlights the reading order for user, developer, and advanced technical paths.
% FIGURE-SPEC-END
\begin{figure}[htbp]
\centering
\figureplaceholder{figs/generated/manual-reading-map}
\caption{Recommended progression through the manual from first contact to advanced topics.}
\label{fig:manual-reading-map}
\end{figure}

Figure~\ref{fig:manual-audience-matrix} maps audiences to the chapters that matter most to them.

% FIGURE-SPEC-BEGIN
% ID: fig:manual-audience-matrix
% Source asset: figs/generated/manual-audience-matrix
% Caption: Audience-to-chapter matrix for the current manual structure.
% Accessibility: A matrix that helps the reader identify the chapters most relevant to a given role or goal.
% FIGURE-SPEC-END
\begin{figure}[htbp]
\centering
\figureplaceholder{figs/generated/manual-audience-matrix}
\caption{Audience-to-chapter matrix for the current manual structure.}
\label{fig:manual-audience-matrix}
\end{figure}

\section*{Safety and Scientific Validity}
\addcontentsline{toc}{section}{Safety and Scientific Validity}
\label{sec:front-matter-safety-and-scientific-validity}

The manual uses a bounded evidence ladder: build evidence, startup evidence, functional evidence, numerical evidence, scientific evidence, security evidence, and release readiness.

Figure~\ref{fig:manual-evidence-ladder} summarizes that ladder and the distinctions between the evidence levels.

% FIGURE-SPEC-BEGIN
% ID: fig:manual-evidence-ladder
% Source asset: figs/generated/manual-evidence-ladder
% Caption: Evidence ladder used by the manual.
% Accessibility: A hierarchical ladder that distinguishes build, startup, functional, numerical, scientific, security, and release evidence.
% FIGURE-SPEC-END
\begin{figure}[htbp]
\centering
\figureplaceholder{figs/generated/manual-evidence-ladder}
\caption{Evidence ladder used by the manual.}
\label{fig:manual-evidence-ladder}
\end{figure}

Figure~\ref{fig:manual-application-scope} summarizes the current application scope and the maturity level claimed for each family of tools.

% FIGURE-SPEC-BEGIN
% ID: fig:manual-application-scope
% Source asset: figs/generated/manual-application-scope
% Caption: Current application scope and maturity snapshot.
% Accessibility: A compact summary of which applications are validated for build or startup and which remain partially validated or planned.
% FIGURE-SPEC-END
\begin{figure}[htbp]
\centering
\figureplaceholder{figs/generated/manual-application-scope}
\caption{Current application scope and maturity snapshot.}
\label{fig:manual-application-scope}
\end{figure}

This book does not treat build success as scientific validation, and local-only or private-only services remain documented with that boundary intact.

\section*{Conventions}
\addcontentsline{toc}{section}{Conventions}
\label{sec:front-matter-conventions}

The canonical prose is written in English. Commands, file names, targets, labels, and code symbols are preserved as they appear in the repository.

The manual uses the following maturity words in a narrow, technical sense:

\begin{itemize}
\item supported;
\item partially validated;
\item experimental;
\item planned.
\end{itemize}

Figure blocks are documented with `FIGURE-SPEC` comments and a placeholder asset until a final image is available. That rule keeps the source explicit even when the artwork is still pending.


This manual remains a draft and should be treated as a work in progress. It is intentionally conservative about claims: build and startup success are reported as build or startup evidence only, and stronger claims require the corresponding validation level. The \textit{Worker} service remains documented with its local or private boundary intact.

The front matter sections are organized in Sections~\ref{sec:front-matter-repository-revision}, \ref{sec:front-matter-reading-paths}, \ref{sec:front-matter-safety-and-scientific-validity}, and \ref{sec:front-matter-conventions}.

\noindent Florianopolis, \datalonga.
